% ============================================================
% PROYECTO INTEGRADOR: PANEL DIDÁCTICO CON MAZO DE CABLES
% ============================================================
\clearpage
\section*{PROYECTO INTEGRADOR}
\subsection*{Panel de instrumentos didáctico con su mazo de cables}

\begin{recuadro}
\textbf{Consigna general:} diseñar y construir un \textbf{panel de instrumentos didáctico} para el
taller, con instrumentos digitales reales o simulados, su mazo de cables completo hecho con
materiales y técnicas aeronáuticas, y toda la documentación de una instalación real.

\textbf{Modalidad:} grupos de cuatro o cinco estudiantes. \textbf{Entrega final:} panel funcionando,
carpeta técnica y presentación oral.
\end{recuadro}

% ------------------------------------------------------------
\section*{1. ALCANCE}
El grupo elige una de estas opciones según el equipamiento disponible:
\begin{itemize}
  \item \textbf{Opción A -- Equipos reales:} panel con un instrumento digital (G5, GI 275 o
        similar), su magnetómetro y, si se dispone, un transpondedor o radio, cableados según el
        manual de instalación del fabricante, con líneas pitot y estática para ensayarlo con el
        banco.
  \item \textbf{Opción B -- Instrumento digital de construcción propia:} un ``EFIS didáctico''
        hecho con un microcontrolador (por ejemplo, ESP32), un sensor de giróscopos y
        acelerómetros, un sensor de presión barométrica y una pantalla, que muestre actitud,
        altitud y velocidad vertical. Se cablea con conectores D-sub y circulares, cables
        aeronáuticos y un disyuntor, como si fuera un equipo certificado.
  \item \textbf{Opción C -- Mazo de instalación completo:} el mazo de cables para la instalación de
        un sistema real (por ejemplo, un G3X Touch con sus sensores), construido a partir de los
        diagramas del manual de instalación, con cajas que simulan cada equipo y sus conectores.
\end{itemize}

% ------------------------------------------------------------
\section*{2. ETAPAS}
\begin{center}
\small
\begin{tabular}{|c|p{0.62\linewidth}|c|}
\hline
\textbf{Etapa} & \textbf{Tareas} & \textbf{Semanas} \\ \hline
1 & Estudio del manual de instalación; lista de materiales con números de parte; análisis de
carga eléctrica. & 2 \\ \hline
2 & Diseño del panel (escala 1:1), plano de cableado y tabla de cables. & 3 \\ \hline
3 & Construcción del panel: recortes, montaje, placas e iluminación. & 2 \\ \hline
4 & Armado del mazo en tablero: crimpado, conectores, blindajes, identificación. & 4 \\ \hline
5 & Verificación pin a pin y verificación de tensiones antes de conectar los equipos. & 1 \\ \hline
6 & Configuración, calibración y ensayos funcionales (y pitot-estático, en la opción A). & 2 \\
\hline
7 & Carpeta técnica, registros e instrucciones de mantenimiento; presentación oral. & 2 \\ \hline
\end{tabular}
\end{center}

% ------------------------------------------------------------
\section*{3. CARPETA TÉCNICA}
\begin{enumerate}
  \item Descripción del sistema y diagrama de bloques.
  \item Lista de materiales con números de parte y certificados de los materiales usados.
  \item Plano del panel y plano de cableado con su tabla de cables.
  \item Registros de inspección de crimpados y ensayos de tracción.
  \item Planilla de verificación pin a pin y de tensiones.
  \item Configuración del equipo y registros de calibración.
  \item Resultados de los ensayos funcionales.
  \item Instrucciones de mantenimiento del panel y tarjeta de trabajo.
\end{enumerate}

% ------------------------------------------------------------
\section*{4. EVALUACIÓN}
\begin{center}
\small
\begin{tabular}{|p{0.66\linewidth}|c|}
\hline
\textbf{Criterio} & \textbf{Peso} \\ \hline
Diseño: plano del panel, plano de cableado y lista de materiales & 20\% \\ \hline
Calidad del mazo: crimpados, conectores, blindajes, atado e identificación & 30\% \\ \hline
Verificaciones y ensayos realizados y documentados & 20\% \\ \hline
Funcionamiento del panel & 15\% \\ \hline
Carpeta técnica y presentación oral & 15\% \\ \hline
\end{tabular}
\end{center}

\begin{recuadro}
\textbf{Para el CV:} \textit{``Diseño y construcción de un panel de instrumentos digitales con su
mazo de cables según prácticas aeronáuticas (AC 43.13-1B): crimpado con herramientas M22520,
conectores MIL-DTL-38999 y D-sub, blindajes, verificación pin a pin, configuración y ensayos''}.
\end{recuadro}
